\documentclass[10pt]{memoir}

% --- Desactivar encabezados y pies ---
\pagestyle{empty}
\setlength{\headheight}{0pt}
\setlength{\headsep}{0pt}
\setlength{\footskip}{0pt}

\setstocksize{3in}{5in}
\settrimmedsize{3in}{5in}{*}
\settypeblocksize{2.8in}{4.8in}{*}
\setlrmargins{0.1in}{*}{*}
\setulmargins{0.1in}{*}{*}
\checkandfixthelayout

\usepackage[spanish]{babel}
\usepackage[utf8]{inputenc}
\usepackage{amsmath, amssymb}

\begin{document}

    \clearpage
    \begin{center}
        \textbf{\Large ¿Puede haber ciclos en una red neuronal?} \\
    \end{center}
    
    \clearpage
    \begin{center}
        \textbf{\Large ¿Puede haber ciclos en una red neuronal?} \\
    \end{center}
        No. 

    \clearpage
    \begin{center}
        \textbf{\Large Explicar la estructura del perceptrón} \\
    \end{center}

    \clearpage
    \begin{center}
        \textbf{\Large Explicar la estructura del perceptrón} \\
    \end{center}
        Se cuenta con los inputs (escalares), un peso de ese input y un bias (opcionalmente). \\
        Se hace la suma de los inputs multiplicados por sus respectivos pesos. \\
        Finalmente, se aplica una función de activación.

    \clearpage
    \begin{center}
        \textbf{\Large ¿Qué es un MLP? ¿Cómo se organiza?} \\
    \end{center}

    \clearpage
    \begin{center}
        \textbf{\Large ¿Qué es un MLP? ¿Cómo se organiza?} \\
    \end{center}
        Un MLP es un perceptrón multicapa. \\
        Un MLP cuenta con todas capas fully connected (FC). Tiene una input layer, hidden layer(s) y una output layer.

    \clearpage
    \begin{center}
        \textbf{\Large Definir shallow y deep network} \\
    \end{center}

    \clearpage
    \begin{center}
        \textbf{\Large Definir shallow y deep network} \\
    \end{center}
        Una shallow network tiene más neuronas de input. Las hace más anchas viéndolas de manera horizontal. \\
        Una deep network tiene más hidden layers. La hace más profunda.

    \clearpage
    \begin{center}
        \textbf{\Large ¿Qué es un tensor? ¿Qué ventajas tiene?} \\
    \end{center}

    \clearpage
    \begin{center}
        \textbf{\Large ¿Qué es un tensor? ¿Qué ventajas tiene?} \\
    \end{center}
        Un tensor es un arreglo multidimensional de escalares. \\
        Sirve para datos de alta dimensionalidad (como imágenes o audios), es eficiente para GPU, y es una representación fácil de inputs, outputs, weights y biases.

    \clearpage
    \begin{center}
        \textbf{\Large Definir forward propagation y sus pasos a seguir.} \\
    \end{center}

    \clearpage
    \begin{center}
        \textbf{\Large Definir forward propagation y sus pasos a seguir.} \\
    \end{center}
        Forward propagation computa el output para un input. \\
        \\
        Se hacen los cálculos a través de cada capa en el sentido de IL/HL/OL para conseguir una predicción $\hat{y}$. A medida que se avanza, se guardan los parámetros para usar en backward propagation.

    \clearpage
    \begin{center}
        \textbf{\Large Definir backward propagation y su proceso. ¿Qué ventaja y desventaja tiene?} \\
    \end{center}

    \clearpage
    \begin{center}
        \textbf{\Large Definir backward propagation y su proceso. ¿Qué ventaja y desventaja tiene?} \\
    \end{center}
        Es el proceso de actualizar los pesos y biases usando el gradiente de la función loss. Se busca minimizar la función loss en el training set. Va en sentido contrario a forward propagation. \\
        \\
        Primero, se calcula la función loss definida por $\hat{y}-y$ a grandes rasgos. \\
        Luego, se calculan los gradientes usando la regla de la cadena. Luego, se aplica el learning rate \textit{LR} al gradiente.
        Finalmente, se actualizan los pesos y biases yendo en sentido contrario: OL/HL/IL. \\
        \\
        Ventaja: es eficiente, ya que una observación del dataset requiere una pasada forward y una pasada backward. \\
        Desventaja: es sensible a la arquitectura de la red, porque afecta su estabilidad.

    \clearpage
    \begin{center}
        \textbf{\Large Nombrar los dos tipos de capas.} \\
    \end{center}

    \clearpage
    \begin{center}
        \textbf{\Large Nombrar los dos tipos de capas.} \\
    \end{center}
        Capas entrenables: sus pesos y biases se modifican en backpropagation. \\
        Capas no entrenables: sus pesos y biases quedan fijos en training para mantener las features aprendidas previamente.

    \clearpage
    \begin{center}
        \textbf{\Large ¿Qué es la vectorización y para qué sirve?} \\
    \end{center}

    \clearpage
    \begin{center}
        \textbf{\Large ¿Qué es la vectorización y para qué sirve? ¿Cuáles son sus ventajas? ¿Cómo queda el cálculo de una neurona en una hidden layer?} \\
    \end{center}
        Es una práctica común para almacenar pesos, inputs y biases en vectores. \\
        Sirve para simplificar las operaciones y acelerar los tiempos. \\
        \\
        Las ventajas que trae son la escalabilidad, la claridad, la capacidad de procesamiento en batch, performance y paralelismo. \\
        \\
        \[\phi(\mathbf{W}^{\top}\mathbf{x}+\mathbf{b})\]
        Donde $\phi$ es la función de activación, $\mathbf{W}$ el vector de pesos, $\mathbf{x}$ el vector de inputs y $\mathbf{b}$ el vector de biases.

    \clearpage
    \begin{center}
        \textbf{\Large ¿Qué es y para qué sirve la función loss?} \\
    \end{center}

    \clearpage
    \begin{center}
        \textbf{\Large ¿Qué es y para qué sirve la función loss?} \\
    \end{center}
        Es una función que compara la predicción predicha con la etiqueta real en el proceso de training: \(\hat{y}-y\). \\
        Sirve para guiar el entrenamiento de la red.

    \clearpage
    \begin{center}
        \textbf{\Large ¿Para qué sirven los grafos computacionales?} \\
    \end{center}

    \clearpage
    \begin{center}
        \textbf{\Large ¿Para qué sirven los grafos computacionales?} \\
    \end{center}
    Para reutilizar cálculos y ayudar con la eficiencia del backpropagation. Las librerías lo usan mucho. 

    \clearpage
    \begin{center}
        \textbf{\Large ¿Para qué sirven y por qué son importantes las funciones de activación?} \\
    \end{center}

    \clearpage
    \begin{center}
        \textbf{\Large ¿Para qué sirven y por qué son importantes las funciones de activación?} \\
    \end{center}
        Sirven para aprender mapeos no lineales y para capturar estructuras complejas de los datos. \\
        \\
        Porque, de no tenerlas, la red se colapsaría a una única capa con la combinación lineal.

    \clearpage
    \begin{center}
        \textbf{\Large Nombrar y describir las funciones de activación lineales. ¿Qué ventajas y qué desventajas tiene?} \\
    \end{center}

    \clearpage
    \begin{center}
        \textbf{\Large Nombrar y describir la función de activación lineal. ¿Qué ventajas y qué desventajas tiene?} \\
    \end{center}
        Función de identidad: el output es el mismo que el input. \(f(x) = x\). \\
        \\
        Ventajas: simples y eficientes, son fáciles de interpretar. \\
        Desventajas: limitadas con datos con patrones complejos, colapsa la red. Generalmente se usa en la capa de output.

    \clearpage
    \begin{center}
        \textbf{\Large Nombrar y describir las funciones escalonadas. ¿Qué ventajas y qué desventajas tienen?}
    \end{center}

    \clearpage
    \begin{center}
        \textbf{\Large Nombrar y describir las funciones escalonadas. ¿Qué ventajas y qué desventajas tienen?} \\
    \end{center}
        Heaviside step: cuando el valor está del lado izquierdo a un umbral tiene un valor. Cuando está del lado derecho tiene otro. \\
        Función de signo: es una función heaviside step pero con -1 para valores negativos y 1 para valores positivos. \\
        \\
        Ventajas: son simples y fáciles de implementar y las salidas binarias convienen para problemas de clasificación. \\
        Desventajas: no es diferenciable en el threshold, es sensible a pequeños cambios en el input y no posee información de gradiente para aprender.

    \clearpage
    \begin{center}
        \textbf{\Large Nombrar y describir las funciones lineales por partes. ¿Qué ventajas y qué desventajas tienen?}
    \end{center}

    \clearpage
    \begin{center}
        \textbf{\Large Nombrar y describir las funciones lineales por partes. ¿Qué ventajas y qué desventajas tienen?} \\
    \end{center}
    ReLU: si el input es negativo, el output es cero. Si el input es positivo, es la función identidad. \(f(x) = \max (0, x)\) \\
    Leaky ReLU: el resultado para los valores negativos cambia. Es igual a ReLU solo que la parte negativa está inclinada. \\
    Shifted ReLU: es ReLU pero con otro threshold. \\
    \\
    Ventajas (ReLU): simple y eficiente, introduce no linealidad en la red, ayuda con el problema de vanishing gradient. \\
    Desventajas (ReLU): las neuronas con el input negativo pueden desactivarse. \\
    \\
    Ventajas (var. ReLU): threshold variable, la neurona no puede morir, mismos beneficios que ReLU. \\
    Desventajas (var. ReLU): más hiperparámetros.

    \clearpage
    \begin{center}
        \textbf{\Large Nombrar y describir las funciones suaves. ¿Qué ventajas y qué desventajas tienen?}
    \end{center}

    \clearpage
    \begin{center}
        \textbf{\Large Nombrar y describir las funciones suaves. ¿Qué ventajas y qué desventajas tienen?} \\
    \end{center}
    Softplus: es la versión suave de la ReLU. \(f(x) = \ln(1+e^{x})\). \\
    ELU: es la versión suave de la Shifted ReLU. \\
    Sigmoide: es una curva con forma de S, con su output acotado entre 0 y 1. \\
    Tanh: es igual que la sigmoide, con su output acotado entre -1 y 1. \\
    Swish: es una combinación entre ReLU y Sigmoide. \\
    \\
    Ventajas: son suaves y diferenciables. La sigmoide sirve para clasificación, tanh para convergencia, swish usa self-gating y ELU evita que la neurona muera. \\
    Desventajas: son costosas y sufren del vanishing gradient.
  
    \clearpage
    \begin{center}
        \textbf{\Large Describir el problema de Vanishing Gradient y Dying Neuron.}
    \end{center}

    \clearpage
    \begin{center}
        \textbf{\Large Describir el problema de Vanishing Gradient y Dying Neuron.}
    \end{center}
    Vanishing Gradient: cuando los gradientes se hacen muy pequeños durante backpropagation a través de las capas. \\
    Dying Neuron: cuando las neuronas siempre están inactivas y sus pesos no se actualizan dado que su gradiente es cero.

    \clearpage
    \begin{center}
        \textbf{\Large Explicar la función Softmax y Log-Softmax.}
    \end{center}

    \clearpage
    \begin{center}
        \textbf{\Large Explicar la función Softmax y Log-Softmax.}
    \end{center}
    Softmax: es una operación que se aplica únicamente y en simultáneo a todas las neuronas en la capa de output. Sirve para normalizar las salidas a una distribución de probabilidad. \\
    Log-Softmax: es hacer el logaritmo de la función SoftMax. Sirve para aportar estabilidad a los números.

    \clearpage
    \begin{center}
        \textbf{\Large }
    \end{center}

    \clearpage
    \begin{center}
        \textbf{\Large }
    \end{center}

    \clearpage
    \begin{center}
        \textbf{\Large }
    \end{center}

    \clearpage
    \begin{center}
        \textbf{\Large }
    \end{center}

    \clearpage
    \begin{center}
        \textbf{\Large }
    \end{center}


    \clearpage
    \begin{center}
        \textbf{\Large }
    \end{center}

    \clearpage
    \begin{center}
        \textbf{\Large }
    \end{center}


    \clearpage
    \begin{center}
        \textbf{\Large }
    \end{center}

    \clearpage
    \begin{center}
        \textbf{\Large }
    \end{center}


    \clearpage
    \begin{center}
        \textbf{\Large }
    \end{center}

    \clearpage
    \begin{center}
        \textbf{\Large }
    \end{center}


    \clearpage
    \begin{center}
        \textbf{\Large }
    \end{center}

    \clearpage
    \begin{center}
        \textbf{\Large }
    \end{center}

    \clearpage
    \begin{center}
        \textbf{\Large }
    \end{center}

    \clearpage
    \begin{center}
        \textbf{\Large }
    \end{center}


    \clearpage
    \begin{center}
        \textbf{\Large }
    \end{center}

    \clearpage
    \begin{center}
        \textbf{\Large }
    \end{center}

    \clearpage
    \begin{center}
        \textbf{\Large }
    \end{center}

    \clearpage
    \begin{center}
        \textbf{\Large }
    \end{center}

    \clearpage
    \begin{center}
        \textbf{\Large }
    \end{center}

    \clearpage
    \begin{center}
        \textbf{\Large }
    \end{center}


    \clearpage
    \begin{center}
        \textbf{\Large }
    \end{center}

    \clearpage
    \begin{center}
        \textbf{\Large }
    \end{center}


    \clearpage
    \begin{center}
        \textbf{\Large }
    \end{center}


\end{document}
